\section{Datasets}
\label{app:data}

We evaluate on three academic domain corpora and the five S2AND author name disambiguation datasets.
Physics comes from the American Physical Society (APS) corpus.
Economics and Psychology are subsets cut out of OpenAlex (a bulk snapshot from March $2026$) based on journals.
We convert every paper into a single string \texttt{Title: \ldots{} Abstract: \ldots} and embed it once.

\textbf{Physics (APS).}
The APS corpus contains $644{,}022$ papers from $18$ \emph{Physical Review} journals ($1893$--$2020$), $8.3$M citation edges, and author-specified PACS topic codes.
Although the public release omits abstracts, we recover them from OpenAlex via DOI matching.
$99.3\%$ of the APS papers match OpenAlex works, and $97.3\%$ of those contain abstracts.
We embed all $644{,}022$ papers regardless of whether we find an abstract.
Every encoder receives the same string.
Missing abstracts therefore introduce no bias into the comparison.
To label clusters, we use PACS nodes from all three levels of the hierarchy (\secref{cluster-labeling}).
Condensed matter accounts for $42.5\%$ of the corpus.

\textbf{Economics and Psychology (OpenAlex).}
For each discipline, we extract Q1/Q2 journals from the Scimago Journal Rank (SJR) in the target subject areas and match them by ISSN to OpenAlex venues.
Economics uses ``Economics, Econometrics and Finance'' ($1{,}035$ journals), and Psychology uses ``Psychology'' ($1{,}166$ journals).
A paper is included if it appears in a matched venue, has a non-empty scientific-work type and title, and has at least one author with $\ge 2$ papers within the slice.
This final condition excludes one-off author--venue combinations and preserves the connectivity of the authorship network.
For text-based tasks, we also require abstracts of $\ge 300$ characters and primary Scopus/ASJC topic labels.
This filter leaves $565{,}475$ Economics papers and $986{,}800$ Psychology papers.
Economics spans $26$ major fields and $244$ subfields.
Psychology spans $26$ major fields and $245$ subfields.
The Scopus/ASJC labels come from OpenAlex and are independent of the SJR venue rules.

% Overlap audit: exps/2026-06-10-general-adapter/compute_leakage_overlap.py -> leakage_overlap.json (issue #22).
% Temporal hardening (issue #72): workflow/rules/groupc_bench.smk, gcb_train_transform.py --cutoff 2018 and the
% temporal topic split in gcb_pool_scores.py; numbers in figs/groupc_temporal.tex. Own pool dir, so Table 1 is untouched.
\textbf{Train-evaluation overlap.}
Both the transformation training pool and the evaluation corpus were obtained from OpenAlex. 
We validate the overlap by checking that evaluation papers account for less than $1\%$ of the training set, and training citation edges falling within a single evaluation field constitute no more than $1.24\%$. Since we do not use any author labels for training $g_\theta$, the author-merging supervision is completely separated. To verify the absence of information leakage, we retrained $g_\theta$ using only citations from $2018$ or earlier; the change in next-paper AUC was at most $.005$, and the ranking of our method remained unaffected. Similarly, even when topic classification is split by time (training $\leq$ 2016, testing $>$ 2016), the encoder's rankings are preserved. For Physics, since the labels are journal codes and the benchmark ended in $2015$, we did not perform a temporal split.

\textbf{Author-name disambiguation (S2AND).}
The five subsets of S2AND~\citep{subramanian2021s2and} are zbMATH ($15{,}181$ papers), QIAN ($6{,}542$), ArnetMiner ($7{,}144$), PubMed ($2{,}871$), and KISTI ($40{,}383$).
Each provides ground-truth author clusters for ambiguous name blocks.
We use title and abstract texts where available, leaving name blocks and metadata untouched.
